\section{Related Work}
\label{sec:related_work}

\smallskip\noindent\textbf{Workload control.}
Our admission controller builds on work that adapts per-input computation at run time under latency, resource, or thermal constraints~\cite{liu-imprecise-computation,shih-imprecise-siam,zilberstein-anytime,loop-perforation-fse,apnet-rtss,branchynet-icpr,msdnet-iclr,mcdnn-mobisys,nestdnn-mobicom,reducto-sigcomm}.
ApproxDet and Chanakya select execution configurations according to input content, resource contention, and latency--accuracy objectives~\cite{approxdet-sensys,chanakya-neurips}.
Recent mobile and embedded systems similarly adapt execution under device constraints: Agile3D selects 3D-detection branches under changing scenes and GPU contention~\cite{agile3d-mobisys}, ARIA coordinates vision inference across heterogeneous mobile processors~\cite{aria-mobisys}, and Phoenix adjusts processor placement and DNN exit depth under thermal throttling~\cite{phoenix-tecs}.
Timecard is closer to our admission logic, tracking each request's elapsed end-to-end budget and adapting its remaining processing to the time left~\cite{timecard-sosp}.
Our admission controller uses a similar budget abstraction, but applies it at stage boundaries, admitting an optional stage only when the remaining draft sequence is still feasible within the capture deadline.
Across this line of work, the control target is the computation performed for the current input.
In our capture pipeline, however, parallel capture can introduce the next capture before the shared draft sequence worker becomes idle.
Reducing current-capture work alone is therefore insufficient.
The pipeline must also control when the next capture is released.

\smallskip\noindent\textbf{Arrival control.}
Our pacing mechanism is related to systems that control request admission or release timing to limit queueing delay~\cite{cruz-network-calculus}.
Networking limits queue growth through active queue management and pacing~\cite{red-floyd-jacobson,codel-cacm,rfc8289,rfc8033,rfc8290,bufferbloat-cacm,tcp-pacing-infocom,bbr-queue}, while server systems admit, reject, or shed requests according to load or deadline signals~\cite{seda-sosp,dagor-socc,breakwater-osdi,aurora-load-shedding,deadline-aware-d3}.
Real-time schedulers similarly regulate job releases by stretching task periods or skipping job instances~\cite{buttazzo-elastic-task,koren-skip-over}.
Preemptive systems can instead interrupt running work to serve short or urgent requests~\cite{shinjuku-nsdi}.
Across these systems, the controller can alter job admission, release, or execution.
In our camera system, a user's shutter press cannot be rejected, and a running stage cannot be preempted.
Our pacing therefore delays next-capture release.
It limits queue growth without reducing work admitted for the current capture.

\smallskip\noindent\textbf{Coordinating workload and arrivals.}
Prior systems have also combined the two control dimensions discussed above.
Feedback-control systems regulate incoming load through request admission while adjusting the service level of admitted requests from deadline-miss or response-time feedback~\cite{lu-feedback-control-edf,lu-feedback-control-scheduling,welsh-overload-control}.
QoS negotiation similarly reduces service quality to accept requests that might otherwise be rejected~\cite{abdelzaher-qos-negotiation}.
These systems control both whether incoming work is accepted and how much service it receives, with rejection remaining one of their control actions.
Content adaptation and brownout instead reduce per-request work without controlling when requests arrive~\cite{abdelzaher-content-adaptation,klein-brownout}.
Neither pattern directly matches our capture pipeline: once issued, a user-triggered capture must be served rather than rejected, and reducing current-capture work alone is insufficient under parallel capture.
Our controller therefore couples the two controls across consecutive captures: admission decides which optional stages the current capture executes, while pacing delays the release of the next capture instead of rejecting it.
To our knowledge, this is the first work to combine per-capture remaining-budget admission with next-capture pacing in a production camera system under a hard per-capture deadline.